\documentclass[10pt,a4paper]{article}
\usepackage{fontspec,amsmath,amssymb,graphicx}
\usepackage[margin=18mm]{geometry}
\setmainfont{Times New Roman}
\setlength{\parindent}{0pt}
\setlength{\parskip}{5pt}
\allowdisplaybreaks
\sloppy
\emergencystretch=2em
\begin{document}
\section*{BT\_4\_5\_3\_03 — Lời giải bổ sung}
\subsection*{Kiến thức và trình tự giải}
\textbf{Bước 1 — Loại đạo hàm cấp một bằng tính đạo hàm đầy đủ.} Với $u_t=Ku_{xx}+Au_x+Bu$, đặt $u=e^{-hx}v$, $h=A/(2K)$. Khi đó


\[
u_t=e^{-hx}v_t,\quad u_x=e^{-hx}(v_x-hv),\quad
u_{xx}=e^{-hx}(v_{xx}-2hv_x+h^2v).
\]


Thay vào PDE, hệ số $v_x$ là $A-2Kh=0$, hệ số $v$ là $Kh^2-Ah+B=B-A^2/(4K)$. Vì thế $v_t=Kv_{xx}+Cv$, $C=B-A^2/(4K)$; có thể tiếp tục đặt $v=e^{Ct}z$ để $z_t=Kz_{xx}$.

\textbf{Bước 2 — Đổi dữ liệu và biên cùng phép đổi ẩn.} $v(x,0)=e^{hx}f(x)$. Dirichlet thuần nhất giữ nguyên; Neumann biến thành $v_x-hv=0$, nên không giữ nguyên một cơ sở cosine nếu $A\ne0$. Với Dirichlet:


\[
u=e^{-hx}\sum_{n\ge1}\left(\frac2L\int_0^Le^{hy}f(y)\sin\frac{n\pi y}Ldy\right)
e^{[C-K(n\pi/L)^2]t}\sin\frac{n\pi x}L.
\]


\textbf{Bước 3 — Tính hệ số mũ–đa thức.} Dùng nguyên hàm


\[
\int e^{hy}\sin ky\,dy
=\frac{e^{hy}(h\sin ky-k\cos ky)}{h^2+k^2}.
\]


Với $y e^{hy}$ hoặc $y(1-y)e^{hy}$, lấy đạo hàm theo $h$ của tích phân đã tính trên cận cố định; trình bày đạo hàm để tìm hệ số thực. Cách này tránh để lại một tích phân sơ cấp chưa thực hiện.

\textbf{Bước 4 — Hệ số không gian biến thiên.} Với $x^2u_{xx}$ trên $x>0$, đặt $z=\ln x$: nếu $U(z,t)=u(e^z,t)$ thì $u_x=U_z/x$, $x^2u_{xx}=U_{zz}-U_z$. Đổi miền và biên sang biến $z$ rồi dùng bước 1. Biên không thuần nhất nên tìm trạng thái dừng trước khi khai triển.

\textbf{Bước 5 — Duy nhất và ổn định.} Với hiệu nghiệm Dirichlet, $I'=-K\int w_x^2+2BI\le2BI$ suy $w=0$ khi $I(0)=0$. Giới hạn phụ thuộc các số $C-K\lambda_n$ có hệ số khác không: âm thì giảm, bằng không thì tồn tại mode dừng, dương thì tăng. Phải kiểm tra tính thống nhất giữa miền PDE và vị trí biên trước khi dùng công thức.
\subsection*{Lời giải áp dụng cho bài}
\textbf{Bước 1. Chỉ ra chỗ thiếu của đề gốc.}
PDE và dữ liệu đầu ghi $0<x<2$, nhưng hai biên đặt tại $0$ và $1$, không có biên tại $2$. Không thể dùng sine $n\pi x$ như một cơ sở Dirichlet trên toàn $(0,2)$ rồi tự coi $1$ là đầu phải. Đây không phải bài Dirichlet trên đoạn được ghi. Trong lớp nghiệm cổ điển với $t>0$ và liên tục đến dữ liệu đầu trong miền, còn có mâu thuẫn: nghiệm trên $(0,1)$ được xác định duy nhất bởi hai biên $0$. Nghiệm nhiệt hệ số hằng giải tích theo $x$ ở thời gian dương, nên thác triển qua điểm trong $1$ buộc trùng phản xạ lẻ của nghiệm sine. Khi $1<x<2$, mặt đầu thác triển là $-f(2-x)=x^2-3x+2$, khác $f=x-x^2$ đã in, với hiệu $-2(x-1)^2$. Vì vậy cần sửa miền hoặc biên/dữ liệu để có nghiệm trong lớp này.
\textbf{Bước 2. Nếu miền dự kiến là $(0,1)$.}
Đây là giả thiết sửa đề, cần ghi riêng. Khi ấy $X_n=\sin n\pi x$, $\lambda_n=n^2\pi^2$. Với $f=x-x^2$,
\[c_n=2\int_0^1f\sin kx\,dx
=-\frac2{k^2}\int_0^1(-2)\sin kx\,dx
=\frac{4(1-(-1)^n)}{k^3},\quad k=n\pi.\]
Mode thỏa $T_n'=(\alpha-k^2)T_n$:
\[u=\sum_{\substack{n\ge1\\n\ {\rm lẻ}}}\frac8{n^3\pi^3}
e^{(\alpha-n^2\pi^2)t}\sin n\pi x,\quad 0<x<1.\]
Sine bằng $0$ ở $0,1$, chuỗi đầu bằng $x-x^2$, và $u_t-u_{xx}=\alpha u$ theo từng mode. Nếu $\alpha<\pi^2$, giới hạn $0$; nếu bằng $\pi^2$, giới hạn $8\sin\pi x/\pi^3$; nếu lớn hơn, mode đầu tăng. Những kết quả này giải bài đã sửa miền, không tự bổ sung điều kiện cho phần $(1,2)$ của đề gốc.\end{document}